\documentclass[12pt,a4paper]{article}
\usepackage{amsmath,amssymb,bm,graphicx,geometry}
\usepackage{physics}
\usepackage{caption}
\usepackage{siunitx}
\geometry{margin=1in}
\setlength{\parskip}{0.8em}
\setlength{\parindent}{0pt}

\begin{document}

\title{\textbf{Jacobian and Singularity Analysis for the Mobile Manipulator}}
\author{Reference Notes -- Muhammed Elyamani}
\date{}
\maketitle

\section*{1. Background}

For a robotic system with generalized coordinates
\[
q =
\begin{bmatrix}
x_m & y_m & \theta_m & \theta_r & \theta_l &
\theta_1 & \theta_2 & \theta_3
\end{bmatrix}^{T},
\]
the velocity of a task point (e.g., the center of mass of link~3) is
\[
\dot{\bm{p}} = J(q)\,\dot{q},
\]
where \(J(q)\in\mathbb{R}^{3\times 8}\) is the \emph{geometric Jacobian} that maps joint and base velocities to Cartesian velocity of the point.

---

\section*{2. Full and Reduced Jacobians}

\subsection*{2.1 Full Jacobian}
The \texttt{full} Jacobian $J_{\text{full}}$ is obtained by differentiating the COM position with respect to \emph{all} generalized coordinates:
\[
J_{\text{full}} = \frac{\partial \bm{p}}{\partial q}.
\]
It represents the kinematic coupling of every coordinate---including nonholonomic base variables---to the end--effector motion.
The scalar manipulability measure is often computed as
\[
\mu_{\text{full}} = \det(J_{\text{full}}J_{\text{full}}^{T}).
\]
Although informative geometrically, $\mu_{\text{full}}$ overestimates the dexterity of the real robot because
$x_m$, $y_m$, and $\theta_m$ are \emph{not independently actuated}.

\vspace{1em}
\textbf{Example (symbolic result)}:
\begin{align*}
\mu_{\text{full}}
&= 0.0051\cos\theta_1 + 0.0229\cos\theta_2
   + 0.0149\cos\theta_1\cos\theta_2
   + \cdots + 0.0564,
\end{align*}
a long trigonometric expression capturing global coupling of all joints.

---

\subsection*{2.2 Differential--Wheel + Arm Jacobian}

In practice, only the wheel angles $(\theta_r,\theta_l)$ and the arm joints $(\theta_1,\theta_2,\theta_3)$ are directly actuated.  
The corresponding Jacobian subset is
\[
J_{\text{mix}} =
\begin{bmatrix}
J_{\text{full}}(:,\, [\theta_r,\theta_l,\theta_1,\theta_2,\theta_3])
\end{bmatrix},
\]
leading to the manipulability measure
\[
\mu_{\text{mix}} = \det(J_{\text{mix}}J_{\text{mix}}^{T}).
\]

After parameter substitution and symbolic simplification,
\[
\boxed{
\mu_{\text{mix}}
= 5.0625\times10^{-6}
\,\sin^{2}\theta_3
\,F(\theta_2,\theta_3)
}
\]
with
\begin{align*}
F(\theta_2,\theta_3) &=
18\cos(\theta_3 + 2\theta_2)
+18\cos(2\theta_2)
+4.5\cos(2\theta_3 + 2\theta_2)
\\ &\quad
+24\cos(\theta_3 + \theta_2)
+48\cos(\theta_2)
+18\cos(\theta_3)
+38.5.
\end{align*}

---

\section*{3. Interpretation}

\begin{itemize}
\item $\mu_{\text{mix}} \ge 0$ always, since $J J^{T}$ is positive semidefinite.
\item Singularities occur when $\mu_{\text{mix}} = 0$:
\[
\sin\theta_3 = 0 \quad \text{or} \quad F(\theta_2,\theta_3)=0.
\]
\end{itemize}

\subsection*{3.1 Primary (Elbow) Singularity}
\[
\sin\theta_3 = 0 \quad\Rightarrow\quad \theta_3 = k\pi.
\]
The forearm and upper arm align (fully stretched or folded),
causing loss of one independent direction of motion.
This is the classic ``elbow singularity'' that eliminates vertical or radial motion of the COM.

\subsection*{3.2 Coupling Singularity}
When \(F(\theta_2,\theta_3)=0\),
certain combinations of shoulder and elbow elevation make
the wheel--induced motion directions line up with the arm's motion, again
reducing the rank of $J_{\text{mix}}$.
These correspond to specific \((\theta_2,\theta_3)\) surfaces in configuration space.

---

\section*{4. Practical Implications}

This section summarizes how the Jacobian and its singular values influence
\emph{real} motion control, trajectory planning, and numerical stability
for the mobile manipulator.  
While the previous derivations identify the singular conditions analytically,
the following points describe how to interpret and use them in practice.

\subsection*{4.1 Most Sensitive Joints}
The analysis of $\mu_{\text{mix}}$ shows that the manipulator's
configuration is dominated by the shoulder and elbow joints
$(\theta_2,\theta_3)$, which primarily control the arm's
\emph{elevation} and \emph{reach}.
When these joints become aligned (for instance,
$\theta_3 \approx 0$ or $\pi$), the arm loses one independent direction of
motion and the Jacobian columns become nearly dependent.
In such cases:
\begin{itemize}
  \item The robot can still move, but only along a restricted set of directions.
  \item Small end--effector motions may require very large joint velocities.
  \item Control signals may saturate or oscillate due to numerical
        amplification of errors.
\end{itemize}
Therefore, avoiding extended or folded elbow configurations is essential
for maintaining full dexterity.

\subsection*{4.2 Dexterity and Conditioning Metrics}
A reliable measure of how ``healthy'' a configuration is
comes from the \textbf{smallest singular value} of the Jacobian,
denoted $\sigma_{\min}$.
If
\[
J = U\,\Sigma\,V^{T}, \qquad
\Sigma = \mathrm{diag}(\sigma_1,\sigma_2,\sigma_3),
\]
then $\sigma_{\min}$ represents the smallest stretching factor of the
mapping from joint velocity space to task space.
\begin{itemize}
  \item Large $\sigma_{\min}$ $\Rightarrow$ well--conditioned;
        the end--effector can move freely in all directions.
  \item Small $\sigma_{\min}$ $\Rightarrow$ near a singular posture;
        some task directions become difficult or impossible to achieve.
  \item $\sigma_{\min} = 0$ $\Rightarrow$ true singularity (rank deficiency).
\end{itemize}
In simulations or experiments, monitor either $\sigma_{\min}$
or the condition number
$\kappa = \sigma_{\max}/\sigma_{\min}$.
A typical warning threshold is
$\sigma_{\min} < 10^{-3}$ or $\kappa > 10^{3}$,
depending on system scaling.

\subsection*{4.3 Exploiting Redundancy and Base Motion}
When the manipulator approaches a low--dexterity region,
the extra degrees of freedom of the system
can be used to ``escape'' singularity.
Two effective strategies are:
\begin{enumerate}
  \item Adjust the azimuth joint $\theta_1$
        to reorient the arm around the vertical axis,
        restoring independent directions of motion.
  \item Move the mobile base slightly using the wheel actuators
        to modify the global pose of the arm,
        thereby improving the local conditioning of $J_{\text{mix}}$.
\end{enumerate}
Such redundancy resolution can be formulated through secondary objectives,
for example,
\[
\dot{q} = J^{+}\dot{\bm{p}}_{\mathrm{des}}
          + (I - J^{+}J)\dot{q}_{\mathrm{null}},
\]
where $\dot{q}_{\mathrm{null}}$ steers the robot away from singularities
without affecting the primary end--effector task.

\subsection*{4.4 Visualization and Workspace Mapping}
To build intuition about safe and unsafe regions,
evaluate $\mu_{\text{mix}}(\theta_2,\theta_3)$ or $\sigma_{\min}(\theta_2,\theta_3)$
over a dense grid covering the feasible joint limits.
Plotting these quantities as a color map reveals:
\begin{itemize}
  \item High--manipulability zones (bright colors)
        where the arm can move easily in all directions.
  \item Narrow valleys (dark colors)
        corresponding to singular surfaces in the $(\theta_2,\theta_3)$ plane.
\end{itemize}
This graphical insight is invaluable for trajectory planners,
which can be programmed to stay within high--manipulability regions.

\subsection*{4.5 Controller Design Considerations}
Controllers based on inverse kinematics or feedback linearization
depend directly on $J$ or its pseudoinverse $J^{+}$.
Close to singularities:
\begin{itemize}
  \item The pseudoinverse amplifies measurement noise,
        causing erratic joint velocity commands.
  \item Numerical inversion becomes unstable or undefined.
  \item Force control formulations using $J^{T}$ lose the ability
        to resist loads in certain directions.
\end{itemize}
In practice, implement a \emph{damped least--squares} inverse,
\[
J^{\#} = J^{T}(J J^{T} + \lambda^2 I)^{-1},
\]
where $\lambda$ is a small damping factor (e.g., $10^{-3}$),
which stabilizes the control law in near--singular zones.

\subsection*{4.6 Summary of Practical Guidelines}
\begin{enumerate}
  \item Keep $\theta_3$ away from $0$ and $\pi$ whenever possible.
  \item Monitor $\sigma_{\min}$ or $\det(JJ^{T})$ during motion.
  \item Use redundant DOFs (base or $\theta_1$) to avoid singular surfaces.
  \item Employ damped or weighted pseudoinverses for robust control.
  \item Visualize and log manipulability maps to guide path planning.
\end{enumerate}

\bigskip
\textbf{In summary:}
the elevation joints $(\theta_2,\theta_3)$ govern the robot's
dexterity envelope.
By continuously monitoring their effect on the Jacobian and
avoiding configurations that degrade $\sigma_{\min}$,
the mobile manipulator maintains stable, smooth, and fully controllable motion
throughout its workspace.


---

\section*{5. Summary Table}

\begin{center}
\begin{tabular}{|l|p{9cm}|}
\hline
\textbf{Perspective} & \textbf{Meaning of Singularity} \\ \hline
Kinematic &
Some task--space velocities $\dot{\bm{p}}$ cannot be produced for any joint velocities $\dot{q}$. \\
\hline
Geometric &
Joint axes or links become aligned; a degree of freedom is lost. \\
\hline
Control / Inverse Kinematics &
$J^{-1}$ or $J^{+}$ becomes ill--conditioned, leading to infinite joint velocities. \\
\hline
Force (Dual) &
Certain end--effector forces require infinite torques; the robot becomes mechanically weak. \\
\hline
\end{tabular}
\end{center}

\bigskip
\textbf{Conclusion:}
For the studied mobile manipulator,
singular behavior is driven mainly by the arm’s elevation configuration
(\(\theta_2,\theta_3\)).
Maintaining $\theta_3$ away from $0$ or $\pi$ and avoiding
combinations that nullify $F(\theta_2,\theta_3)$
ensures full rank of the Jacobian and reliable motion control.

\end{document}
